\subsection{Плавление и кристаллизация}
%%%%%%%%%%%%%%%%%%%%%%%%%%%%%%%%%%%%%%%%%%%%%%%%%%%%
%%%%%%%%%%%%%%%%%%%%%%%%%%%%%%%%%%%%%%%%%%%%%%%%%%%%
%%%%%%%%%%%%%%%%%%%%%%%%%%%%%%%%%%%%%%%%%%%%%%%%%%%%
 \z{Алюминиевому бруску массой $m = 500\,\text{г}$ при температуре $t_1 = 30\,\degree\text{C}$ передали количество теплоты $Q = 900\,\text{кДж}$. Какой станет температура алюминия?}
%
{$t_{\text{Al}} = \frac{Q - c_{\text{тв}} m(t_{\text{пл}} - t_0) - \lambda m}{c_{\text{ж}} m} + t_{\text{пл}} = 1415\,\degree\text{C}$}

%%%%%%%%%%%%%%%%%%%%%%%%%%%%%%%%%%%%%%%%%%%%%%%%%%%%
 \z{В калориметре находится вода массой $m_{\text{в}} = 2\,\text{кг}$, температура которой $t_{\text{в}} = 30\,\degree\text{C}$. В калориметр помещают лед при температуре $0\,\degree\text{C}$. Какова могла быть масса $m_{\text{л}}$ льда, если он растаял?}
%
{$m_{\text{л}} < \frac{c_{\text{в}} m_{\text{в}} t_{\text{в}}}{\lambda} = 760\,\text{г}$}

%%%%%%%%%%%%%%%%%%%%%%%%%%%%%%%%%%%%%%%%%%%%%%%%%%%%
 \z{В калориметре находится лед и вода при температуре $0\,\degree\text{C}$. Масса льда и воды одинаковы и равны $500\,\text{г}$. В калориметр наливают воду массой $m = 1\,\text{кг}$ при температуре $t = 50\,\degree\text{C}$. Какая температура установится в калориметре?}
%
{$t = \frac{c_{\text{в}} m_{\text{в1}} t_{\text{н}} - m_{\text{л}} \lambda}{c_{\text{в}} (m_{\text{л}} + m_{\text{в1}} + m_{\text{в2}})} \approx 5\,\degree\text{C}$}

%%%%%%%%%%%%%%%%%%%%%%%%%%%%%%%%%%%%%%%%%%%%%%%%%%%%
 \z{Кусок льда массой $m_{\text{л}} = 700\,\text{г}$ поместили в калориметр с водой. Масса воды $m_{\text{в}} = 2.5\,\text{кг}$, начальная температура $t_{\text{в}} = 5\,\degree\text{C}$. Когда установилось тепловое равновесие, оказалось, что масса воды увеличилась на $\Delta m = 64\,\text{г}$. Определите начальную температуру льда.}
%
{$t_{\text{л}} = \frac{m_{\text{л}} \lambda - c_{\text{в}} m_{\text{в}} t_{\text{н}}}{c_{\text{л}} m_{\text{л}}}\approx -21\,\degree\text{C}$}

%%%%%%%%%%%%%%%%%%%%%%%%%%%%%%%%%%%%%%%%%%%%%%%%%%%%
 \z{В калориметре при температуре $0\,\degree\text{C}$ находится вода массой $m_{\text{в}} = 500\,\text{г}$ и лед массой $m_{\text{л}} = 300\,\text{г}$. Какая температура установится в калориметре, если долить в него $m_1 = 100\,\text{г}$ кипятка? $m_2 = 500\,\text{г}$ кипятка?}
%
{$t = 0\,\degree\text{C},\ \text{т.к.}\ m_{\text{л}} \lambda > c_{\text{в}} m_{\text{в2}} t_{\text{кип}};\quad t = \frac{c_{\text{в}} m_{\text{в2}} t_{\text{кип}} - m_{\text{л}} \lambda}{c_{\text{в}} (m_{\text{л}} + m_{\text{в1}} + m_{\text{в2}})} \approx 20\,\degree\text{C}$}

%%%%%%%%%%%%%%%%%%%%%%%%%%%%%%%%%%%%%%%%%%%%%%%%%%%%
 \z{В калориметр, содержащий воду массой $m_{\text{в}} = 2\,\text{кг}$ при температуре $t_{\text{в}} = 20\,\degree\text{C}$, бросили ледяной шар массой $m_{\text{л}} = 1\,\text{кг}$, в центр которого вморожен стальной шарик массой $m_{\text{ст}} = 50\,\text{г}$ при температуре $0\,\degree\text{C}$. Где окажется стальной шарик после установления теплового равновесия?}
%
{Будет плавать.}

%%%%%%%%%%%%%%%%%%%%%%%%%%%%%%%%%%%%%%%%%%%%%%%%%%%%
 \z{В переохлажденной до температуры $-10\,\degree\text{C}$ воде происходит быстрый процесс кристаллизации. Какая часть воды при этом превращается в лед?}
%
{$\alpha \approx 13\%$}